\chapter{Installation and tools}
\label{ch:tools}

The reference implementation of Forge is written in OCaml and built with
Dune. It provides an interpreter, an interactive REPL, a native compiler and a
handful of developer tools. Native compilation needs Clang, which links the
generated code with a small C runtime, and the Boehm garbage collector
(\texttt{bdw-gc}), which manages the memory of compiled programs.

\section{Installing}

\paragraph{From source} Forge needs OCaml 5.1 or later, Dune, the
\texttt{ounit2} test library, Clang and the Boehm garbage collector
(\texttt{brew install bdw-gc} or \texttt{apt install libgc-dev}).\index{installation}

\begin{shellpanel}
\begin{lstlisting}
$ opam switch create 5.1.1 && eval "$(opam env)"
$ opam install dune ounit2
$ git clone https://github.com/LambdaAK/Forge && cd Forge
$ make
\end{lstlisting}
\end{shellpanel}

\paragraph{With Docker} The repository's \texttt{Dockerfile} builds an
image with the interpreter, the REPL and the compiler, so nothing but Docker
is needed:\index{Docker}

\begin{shellpanel}
\begin{lstlisting}
$ docker build -t forge:local .
$ docker run --rm forge:local run /opt/forge/programs/minimal.forge
3
$ docker run --rm -it forge:local repl
$ docker run --rm -v "$PWD":/work -w /work forge:local compile /work/hello.forge
$ docker run --rm -v "$PWD":/work -w /work forge:local /work/a.out
Hello, World!
\end{lstlisting}
\end{shellpanel}

\section{The tools}

All executables are built under \texttt{\_build/default/bin} and can be run
with \texttt{dune exec}. They look for the standard prelude
(\texttt{prelude/prelude.forge}) and the C runtime
(\texttt{runtime/ls\_runtime.c}) by searching upwards from the current
directory; set \texttt{FORGE\_ROOT} to the checkout when running them from
elsewhere.

\begin{center}
\small
\begin{tabularx}{\textwidth}{@{}l X@{}}
\toprule
\textbf{Executable} & \textbf{Purpose} \\
\midrule
\texttt{interpreter} \textit{file} & Run a program with the interpreter.\index{interpreter} \\
\texttt{repl} [\textit{file}] & Interactive session; an optional file is loaded first.\index{REPL} \\
\texttt{compile\_forge} \textit{file} [\textit{out}] & Compile to a native executable (default \texttt{a.out}); also writes the intermediate \texttt{.mir}, \texttt{.ll} and \texttt{.s} files next to the source.\index{native compiler} \\
\texttt{dump\_min\_ir} \textit{file} & Print the program in the compiler's intermediate representation, Min IR. \\
\texttt{parse\_file} \textit{file} & Print the parse tree. \\
\texttt{forge\_hover} \textit{file prelude line col} & Type at a position, for editor integration (reads the buffer on standard input). \\
\bottomrule
\end{tabularx}
\end{center}

The \texttt{Makefile} wraps the common tasks:

\begin{center}
\small
\begin{tabularx}{\textwidth}{@{}l X@{}}
\toprule
\textbf{Target} & \textbf{Effect} \\
\midrule
\texttt{make} & Build everything. \\
\texttt{make run FILE=}\textit{f} & Run \textit{f} with the interpreter. \\
\texttt{make repl [FILE=}\textit{f}\texttt{]} & Start the REPL. \\
\texttt{make compile-forge FILE=}\textit{f} \texttt{[OUT=}\textit{exe}\texttt{]} & Native compilation. \\
\texttt{make dump-ir FILE=}\textit{f} & Show the Min IR of a program. \\
\texttt{make suite} / \texttt{make compiler-suite} & The main test suite / every test suite. \\
\texttt{make docs} & Build this manual (\texttt{documentation/Forge.pdf}). \\
\bottomrule
\end{tabularx}
\end{center}

\section{The REPL}
\label{sec:repl}

The REPL reads a definition or an expression, terminated by \texttt{;;},
and prints its value and type. Commands start with a colon:
\texttt{:type} \textit{expr} shows a type without evaluating,
\texttt{:env} lists the bindings, \texttt{:history} the inputs so far,
\texttt{:help} the commands and \texttt{:quit} leaves.

\begin{shellpanel}
\begin{lstlisting}
$ dune exec ./bin/repl.exe
λ> let square x = x * x;;
Name : square
function : Int -> Int
λ> square 12;;
144 : Int
λ> :type map
Type: (b -> a) -> List<b> -> List<a>
λ> [x | x <- [1...4], x > 2];;
[3, 4] : List<Int>
\end{lstlisting}
\end{shellpanel}

\section{Diagnostics}
\label{sec:diagnostics}

Errors are printed on standard error as \texttt{error:} followed by the
message, and the program stops with exit status~1; the program's own output
up to that point is kept. Syntax errors give the file and line. Warnings ---
for instance about a \fc|case| that does not cover every value --- are printed
as \texttt{warning:} and do not stop the program.
Appendix~\ref{ch:diagnostics} lists the messages.

\begin{center}
\small
\begin{tabular}{@{}ll@{}}
\toprule
\textbf{Exit status} & \textbf{Meaning} \\
\midrule
0 & success \\
1 & a lexical, syntax, type or run-time error \\
2 & wrong command-line usage \\
\bottomrule
\end{tabular}
\end{center}

\section{Editor support}

The \texttt{vscode-forge} directory contains a Visual Studio Code extension
with syntax highlighting and type-on-hover (it runs \texttt{forge\_hover}).
Build it with \texttt{npm ci --prefix vscode-forge \&\& npm --prefix
vscode-forge run compile}.
